\emph{内部值与梯度匹配。}
各开分区中的端点和触点映射光滑。
在 $\Gamma_K\cap\{i<K\}$ 上，等待终点趋于当前点，$U$ 两侧均趋于 $W$；
由 \eqref{eq:waiting-gradient}，等待侧梯度的极限为
\[
\left(\frac{p(s-h)}{si(s-r)},\ \frac{p}{i(s-r)}\right),
\]
恰等于切换条件下的 $\nabla W$。
在 $\Phi=\Phi_B,s>s_B,i<K$ 的界面上，两个等待分支的终点均趋于 $(s_B,K)$；
由 $C_B(s_B,s_B)=0$，两侧函数值相同，而共同的梯度极限为
\[
\left(\frac{p(s-h)}{sK(s_B-r)},\ \frac{p}{K(s_B-r)}\right).
\]
上述极限随界面点连续变化。
在每个固定内部界面点的充分小邻域内，所用端点感染值及导数公式的分母保持为正，
故这些公式确实给出各侧梯度到该界面的连续延拓。

下面说明局部拼接的可微性。
两类内部界面分别是 $i=I_\Gamma(s)$ 和
$i=\Phi_B-s+h\ln s$ 的光滑图像，且在不安全内部不相交。
固定其中一个点 $x_*=(s_*,k(s_*))$，其中 $k$ 表示对应图像函数。
缩小邻域后，坐标变换
\[
\Xi(y_1,y_2)=(s_*+y_1,k(s_*+y_1)+y_2)
\]
将以原点为中心的小矩形 $Q$ 映到该点的不安全内部邻域，
并把 $y_2=0$ 映到界面。
该变换及其逆均为 $C^1$。
令 $v=U\circ\Xi$，则 $v$ 连续，在 $Q\cap\{y_2>0\}$ 和
$Q\cap\{y_2<0\}$ 内分别为 $C^1$。
由已证的梯度匹配以及
$\nabla v(y)=(\mathrm D\Xi(y))^{\mathsf T}\nabla U(\Xi(y))$，
两侧梯度可连续拼接为 $Q$ 上的向量值函数 $P$。
当 $y_2\ne0$ 时，从 $\varepsilon y$ 到 $y$ 沿同侧线段积分，
再令 $\varepsilon\downarrow0$，得到
\[
v(y)-v(0)=\int_0^1 P(ty)\cdot y\dd t.
\]
当 $y_2=0$ 时，由连续性从任一侧取极限，等式仍成立。
因此
\[
v(y)-v(0)-P(0)\cdot y
=\int_0^1 [P(ty)-P(0)]\cdot y\dd t
=o(\lVert y\rVert).
\]
这证明 $v$ 在原点可微且导数为 $P(0)$。
在每个界面点重复此论证，并利用 $P$ 的连续性，得到局部 $C^1$ 拼接。
变回原坐标，$U$ 在不安全内部为 $C^1$。
切换曲线的下端点位于安全边界，上端点位于容量边界，
均不属于此处的不安全内部论域。
